\chapter{Configuration}
\label{ch:configuration}

Une fois l'assistant de mise en service passé, le travail de
configuration du jour 1 se compose de trois blocs : définir vos
\emph{lieux de stockage}\index{lieux de stockage}, peupler les
\emph{données de référence} (genres, rôles de contributeur, états de
volume, types de nœud de localisation) et ajuster les
\emph{paramètres globaux} (seuil de retard, cadence d'auto-purge,
clés de fournisseurs).

Tout se configure depuis le panneau
\texttt{/admin}\index{panneau admin} — il n'y a pas de fichier de
configuration séparé. Le panneau a cinq onglets :

\begin{itemize}
  \item \textbf{Health} — tableau de bord en lecture seule.
  \item \textbf{Users} — créer/désactiver des utilisateurs, changer
        de rôle.
  \item \textbf{Reference data} — les quatre tables taxonomiques.
  \item \textbf{Trash} — voir et restaurer les éléments
        soft-deleted.
  \item \textbf{System} — paramètres globaux.
\end{itemize}

Seul le rôle \texttt{Admin} voit le panneau ; les utilisateurs
\texttt{Librarian} et \texttt{Anonymous} reçoivent un 403 sur toute
route \texttt{/admin/*}.

\section{Lieux de stockage et hiérarchie}\label{sec:locations}

Les lieux de stockage sont les endroits physiques où vivent vos
volumes\index{volume} : une pièce, une étagère, une rangée, une
boîte numérotée, une case précise d'une armoire. mybibli les
organise comme un \textbf{arbre}\index{arbre des lieux} — chaque lieu
a au plus un parent. La hiérarchie est arbitraire : vous décidez à la
fois des niveaux et des libellés.

Une hiérarchie domestique exploitable peut tenir sur trois niveaux :

\begin{lstlisting}
Salon                (pièce)
 |-- Bibliothèque A   (étagère)
 |    |-- Rangée 1    (rangée)
 |    |-- Rangée 2    (rangée)
 |-- Bibliothèque B   (étagère)
 |    |-- Rangée 1    (rangée)
Chambre              (pièce)
 |-- Table de nuit   (étagère)
\end{lstlisting}

Trois niveaux suffisent à retrouver n'importe quel volume en quelques
secondes sans forcer la hiérarchie à épouser chaque recoin de votre
intérieur. Une liste plate (une grosse « étagère » par étagère
physique) fonctionne tout aussi bien pour des collections de
quelques centaines de volumes.

\subsection*{Éditer l'arbre}

Ouvrez \texttt{/locations}. L'arbre est rendu sous forme de liste
imbriquée avec des actions création / renommage / déplacement /
suppression en ligne. Le glisser-déposer \emph{n'est pas} pris en
charge (mybibli applique une CSP stricte sans JavaScript inline) ;
les déplacements se font en éditant un nœud et en choisissant un
nouveau parent dans un menu déroulant.

\subsection*{Types de nœud de localisation}

Chaque nœud de localisation porte un \textbf{type de
nœud}\index{type de nœud} — « pièce », « étagère », « rangée »,
« boîte », ce que vous voulez. Les types sont gérés depuis
\texttt{/admin > Reference data > Location node types}. Ce sont de
purs libellés ; mybibli n'impose pas qu'une « étagère » soit sous
une « pièce ».

\begin{tip}
Renommer un type de nœud propage la modification à chaque lieu qui
l'utilise (le type est stocké par nom dans la table des lieux).
Supprimer un type est refusé tant qu'au moins un lieu y fait
référence.
\end{tip}

\section{Étiquettes : codes V et codes L}\label{sec:labels}

mybibli imprime deux types d'étiquettes à code-barres :

\begin{itemize}
  \item \textbf{Codes V}\index{code V} — un par
        \emph{volume}\index{volume} (l'exemplaire physique sur
        l'étagère). Format : \texttt{V} suivi d'un numéro
        zéro-paddé (par exemple \texttt{V0042}).
  \item \textbf{Codes L}\index{code L} — un par \emph{lieu de
        stockage}. Format : \texttt{L} suivi d'un numéro zéro-paddé
        (par exemple \texttt{L0007}).
\end{itemize}

Les codes V et L sont scannés avec la même douchette code-barres que
vous utilisez pour les ISBN au moment du catalogage. Le préfixe
(\texttt{V} ou \texttt{L}) indique à mybibli quel lookup effectuer.

\subsection*{Imprimer les étiquettes}

Ouvrez \texttt{/locations} ou la page du volume concerné et cliquez
\textbf{Imprimer l'étiquette}. mybibli rend une page HTML
imprimable avec le code-barres Code-128 et le texte lisible. Imprimez
sur papier ordinaire ou sur planches autocollantes — l'encodage est
suffisamment robuste pour qu'une imprimante jet d'encre bas de gamme
suffise pour les codes V comme les codes L.

\begin{tip}
Imprimez vos étiquettes de codes V \emph{avant} de commencer un
catalogage en série. Une planche pré-imprimée de V0001 à V0500
permet de saisir un autocollant, de le coller au dos du livre et de
scanner d'un même mouvement — sans aller-retour entre l'interface
de catalogage et l'imprimante.
\end{tip}

\section{Matériel : douchette code-barres}\label{sec:scanner-hardware}

mybibli est conçu autour d'une douchette code-barres mais n'en exige
pas strictement une : tous les champs ISBN, code V et code L
acceptent la saisie clavier comme repli. Cela dit, dès que la
collection dépasse une poignée de titres, une douchette USB
bon marché est rentabilisée après la première douzaine de scans.

\subsection*{Quoi acheter}

N'importe quelle douchette USB en \emph{mode clavier HID}\index{mode
HID clavier} dans la gamme 15--30~€ convient. Le mode HID est
ce qui rend la douchette sans pilote : l'OS la voit comme un clavier
ordinaire et les chiffres scannés atterrissent directement dans le
champ qui a le focus. Presque toutes les douchettes USB grand public
sortent d'usine dans ce mode ; cherchez les mots ``HID'',
``émulation clavier'' ou ``plug-and-play'' sur la fiche produit.

La douchette doit prendre en charge les deux symbologies que
mybibli utilise :

\begin{itemize}
  \item \textbf{EAN-13}\index{EAN-13} — le code à 13 chiffres au dos
        des livres, CD, DVD et de la plupart des produits grand
        public. Universellement supporté.
  \item \textbf{Code-128}\index{Code-128} — la symbologie utilisée
        par mybibli pour encoder les codes V et codes L sur les
        étiquettes qu'il imprime (voir section~\ref{sec:labels}).
        Universellement supporté.
\end{itemize}

Vous n'avez \emph{pas} besoin d'un imageur 2D (le type qui lit les
QR codes) — une douchette laser ou LED 1D basique suffit. mybibli
n'utilise pas les QR codes.

\subsection*{Configuration}

Trois réglages comptent côté douchette. La plupart sortent d'usine
correctement configurées ; si vos scans se comportent bizarrement,
consultez le manuel de la douchette ou la planche de codes-barres de
configuration livrée dans la boîte :

\begin{itemize}
  \item \textbf{Terminateur / suffixe :} \texttt{Enter} (parfois
        appelé \texttt{CR} ou \texttt{Carriage Return}). Le
        formulaire de catalogage de mybibli soumet
        automatiquement à la pression d'Entrée, donc ce réglage
        transforme chaque scan en un seul aller-retour. Toute autre
        valeur (\texttt{Tab}, pas de terminateur) ne marchera pas
        sans intervention clavier manuelle.
  \item \textbf{Préfixe :} aucun. Certaines douchettes peuvent
        préfixer une chaîne fixe au début du payload — laissez
        cette option désactivée.
  \item \textbf{Disposition clavier :} \emph{doit} correspondre à
        la disposition de l'hôte sur lequel tourne le navigateur. Si
        votre ordinateur utilise AZERTY mais la douchette est
        configurée en QWERTY, les chiffres sortent en lettres (p.~ex.
        \texttt{1234567890} arrive comme \texttt{\&é"'(-è\_çà}) et
        aucun ISBN ne résout. Le manuel de la douchette contient une
        planche de codes pour changer la disposition.
\end{itemize}

\begin{tip}
Douchettes USB sur NAS Synology : peu importe que la douchette soit
branchée sur l'ordinateur portable depuis lequel vous accédez à
l'interface web de mybibli plutôt que sur le NAS lui-même. La
douchette est un clavier pour le navigateur, pas pour le serveur. Le
même montage marche avec une appli scanner sur téléphone, une
douchette Bluetooth industrielle appairée à l'ordinateur, ou
n'importe quelle douchette USB clé.
\end{tip}

\subsection*{Matériel validé sur le terrain}

Comme mybibli n'a besoin que d'une douchette en émulation clavier et
d'un navigateur, l'appareil qui l'exécute peut être un poste fixe, un
portable ou une tablette. Combinaisons validées sur le terrain :

\begin{itemize}
  \item \textbf{iPad (Safari) + lecteur Eyoyo.} Une tablette au poste
        de catalogage avec une douchette Eyoyo sans fil — pratique pour
        parcourir les rayons pendant un inventaire.
  \item \textbf{Poste de travail Linux + Inateck BCST-36.} Un bureau de
        catalogage fixe avec une douchette Bluetooth Inateck BCST-36.
\end{itemize}

Ce sont des exemples, pas une liste blanche de compatibilité —
n'importe quelle douchette en mode clavier HID (voir \emph{Quoi
acheter} ci-dessus) fonctionne de la même façon.

\subsection*{Quand la douchette n'est pas disponible}

Chaque champ de scan accepte aussi la saisie clavier. Cliquez dans
le champ étiqueté ``Scanner ISBN'' (ou ``Scanner code V'' /
``Scanner code L'' ailleurs), tapez les chiffres, pressez Entrée. Le
formulaire se comporte à l'identique.

\section{Données de référence}\label{sec:reference-data}

Quatre tables taxonomiques pilotent les menus déroulants que vous
voyez pendant le catalogage. Elles vivent sous \texttt{/admin >
Reference data} :

\begin{description}
  \item[Genres]\index{genres} Les genres assignés aux titres.
        mybibli livre un seed de base FR/EN (Fiction, Documentaire,
        Biographie, \dots) — éditez ou étendez pour coller à votre
        collection.
  \item[Rôles de contributeur]\index{rôles de contributeur} Auteur,
        Illustrateur, Traducteur, Éditeur, Narrateur, \dots Utilisés
        sur les pages détail de titre pour attribuer les bonnes
        personnes aux bons rôles.
  \item[États de volume]\index{état de volume} L'état physique d'un
        exemplaire — Neuf, Bon, Usé, Endommagé, \dots Utile pour les
        retours de prêt (marquer un volume Usé quand un emprunteur
        l'a abîmé).
  \item[Types de nœud de localisation]\index{type de nœud} Voir
        section~\ref{sec:locations}.
\end{description}

\subsection*{Sémantique CRUD}

Les quatre tables partagent le même schéma CRUD :

\begin{itemize}
  \item \textbf{Créer} ajoute une ligne. Si le nom entre en collision
        avec une ligne soft-deleted de la même table, la ligne
        existante est réactivée et le résultat affiche
        \emph{Réactivé} au lieu de \emph{Créé}.
  \item \textbf{Renommer} met à jour le nom. Pour les types de nœud
        de localisation, le nouveau nom cascade sur tous les lieux
        qui utilisent ce type.
  \item \textbf{Supprimer} effectue un soft-delete \emph{uniquement
        si la ligne est inutilisée}. Si un titre référence encore
        ce genre / rôle / état / type, la suppression est refusée
        avec un comptage des dépendants.
\end{itemize}

\subsection*{Les données de référence ne sont pas traduites}

Les valeurs des données de référence sont stockées telles quelles —
le seed FR crée un genre appelé « Roman » et c'est ce texte
littéral que verront les utilisateurs anglophones s'ils basculent
l'interface en anglais. Seule l'ossature de l'interface\index{i18n}
(libellés de boutons, navigation, messages d'erreur) est traduite.

\section{Paramètres système}\label{sec:system}

L'onglet \texttt{/admin > System} regroupe tous les paramètres
globaux. Le panneau est divisé en sections, chacune avec son bouton
de sauvegarde :

\subsection*{Prêts}

\begin{itemize}
  \item \textbf{Seuil de retard (jours).} Définit combien de jours
        un prêt peut courir avant d'apparaître comme en retard sur
        le tableau de bord. Défaut : 30.
  \item \textbf{Intervalle d'auto-purge (secondes).} Cadence à
        laquelle le scheduler d'arrière-plan supprime
        définitivement les éléments soft-deleted depuis plus de
        30\,jours. Défaut : 86400\,s (24\,heures).
\end{itemize}

\subsection*{Sessions}

Depuis la 1.11.0, le \textbf{délai d'inactivité de session} se règle
ici, en \emph{heures} (1--720, défaut 4). Il pilote deux choses à la
fois : la durée de survie d'une session connectée sans activité, et
la durée de vie du cookie de session dans le navigateur — une
tablette qui verrouille son écran en plein catalogage ne perd donc
plus sa connexion quand le navigateur jette les cookies de session.
Appliqué à la requête suivante ; les sessions déjà ouvertes ne sont
pas coupées rétroactivement.

\subsection*{Fournisseurs de métadonnées}

Une ligne par fournisseur à clé (Google Books, OMDb, TMDb). Chaque
ligne a un champ clé et un bouton \emph{Effacer}. Sauvegarder une
clé prend effet à la prochaine récupération de métadonnées — pas de
redémarrage, pas de redéploiement.

\begin{note}
Si la même clé de fournisseur est positionnée dans \texttt{.env}
\emph{et} effacée depuis l'interface, le prochain reboot re-migre la
valeur depuis \texttt{.env}. L'intention de l'opérateur au moment du
déploiement l'emporte au boot. Pour rendre l'effacement durable,
retirez aussi la variable de \texttt{.env} (ou de
\texttt{docker-compose.yml}).
\end{note}

La même section porte aussi les \textbf{délais d'attente} de
consultation (depuis 1.7.9) : le délai par fournisseur de métadonnées
(temps maximal d'attente d'un fournisseur avant de passer au suivant
dans la chaîne, 5\,s par défaut) et le délai de la sonde de
joignabilité (10\,s par défaut), tous deux bornés à 1--60\,s. Depuis
la 1.9.0, le délai métadonnées peut en plus être \textbf{surchargé
par fournisseur} : le formulaire liste chaque fournisseur enregistré
avec son propre champ ; laissez un champ vide pour utiliser la valeur
globale. Utile quand un fournisseur (une bibliothèque nationale lente,
par exemple) a besoin de plus de marge sans ralentir toute la chaîne.
Les changements s'appliquent au prochain scan — sans redémarrage.

Depuis la 1.11.0, le même formulaire porte aussi le \textbf{délai du
re-fetch de couvertures en lot} (millisecondes, 0--60000, défaut
1000) : la pause entre chaque titre lors d'une exécution de
\emph{Récupérer les couvertures manquantes} depuis l'onglet Santé.
Un rythme trop rapide fait bloquer toute l'exécution par les
fournisseurs — Google Books répond aux rafales par des erreurs et
aucune couverture ne revient. Le défaut ajoute environ deux minutes
par centaine de titres, négligeable pour une action d'arrière-plan.
L'exécution retente désormais un titre bloqué avec un délai croissant
avant d'y renoncer, et se termine par un résumé sur l'onglet Santé :
couvertures récupérées, erreurs fournisseur, et titres pour lesquels
aucune couverture n'existe nulle part.

Depuis la 1.13.0, l'onglet Santé porte une seconde action en lot à
côté du re-fetch de couvertures : \textbf{Compléter les métadonnées
via BnF}. Elle relance la recherche de métadonnées pour \emph{tous}
les titres catalogués disposant d'un code — pas seulement ceux sans
couverture — afin de remplir les champs alignés UNIMARC introduits
en 1.13.0 sur les titres catalogués avant la mise à jour. Elle ne
fait que combler les manques : une valeur déjà présente, ou modifiée
à la main, n'est jamais écrasée. Les deux actions en lot partagent
le même délai de rythme, le même back-off anti-blocage et le même
verrou d'exécution unique — une seule peut tourner à la fois.

\subsection*{Langue}

Fixe la langue d'interface par défaut pour les visiteurs
\emph{anonymes}. Les utilisateurs authentifiés outrepassent ce
choix via le bouton bascule dans la barre de navigation — leur
choix est stocké sur leur ligne utilisateur.

\section{Utilisateurs}\label{sec:users}

Ouvrez \texttt{/admin > Users} pour gérer les comptes utilisateur.
Chaque ligne montre le nom d'utilisateur, le rôle, le drapeau
actif et la date de dernière activité. Les actions disponibles :

\begin{itemize}
  \item \textbf{Créer un utilisateur} (haut de page) — nom, mot de
        passe, rôle.
  \item \textbf{Éditer} — changer le rôle ou renommer. Ne change
        pas le mot de passe d'un autre utilisateur depuis cette
        vue ; les changements de mot de passe passent par la page
        \texttt{/account} de l'utilisateur lui-même.
  \item \textbf{Désactiver} — soft-delete l'utilisateur. Ses
        sessions sont effacées immédiatement ; il ne peut plus se
        reconnecter.
  \item \textbf{Réactiver} — défait une désactivation (depuis
        l'onglet \texttt{Trash}).
\end{itemize}

Deux garde-fous s'appliquent :

\begin{itemize}
  \item \textbf{L'auto-désactivation est bloquée.} Vous ne pouvez
        pas désactiver le compte avec lequel vous êtes connecté.
  \item \textbf{Garde du dernier admin actif.} Le système maintient
        toujours au moins un admin actif. Désactiver le dernier
        est refusé avec un 409 et un message d'erreur clair.
\end{itemize}

Les trois rôles\index{rôles} (Anonyme, Bibliothécaire, Administrateur)
sont détaillés au chapitre~\ref{ch:roles}.
